\chapter{热力学基本定律}

%\section*{宏观}

热力学研究的对象是由大量微观粒子(分子或其他粒子)组成的宏观物质系统。
与系统发生相互作用的其它物体称为外界。
% 由于宏观测量的粗糙性，系统的宏观描述与大多数原子坐标不直接相关。热力学只关心大量原子坐标的平均后的宏观量。

\begin{definition}
	{热力学系统的分类}{thermodynamical system}
	根据系统与外界相互作用的情况，可以将系统分为：
	\begin{itemize}
		\item 孤立系统(isolated system)：与外界既没有物质交换也没有能量交换；
		\item 封闭系统(closed system)：边界不允许物质通过，可以与外界交换能量；
		\item 开放系统(open system)：边界允许物质通过，也可以与外界交换能量。
	\end{itemize}
\end{definition}

\begin{remark}
	由于物质的普遍联系和相互作用，孤立系统的概念实际上只是一个理想的极限概念。
	% 实际情况是，当系统与外界的相互作用十分微弱，交换的粒子数远小于系统本身的粒子数、相互作用的能量远小于系统本身的能量，我们就把系统看作孤立系统。
	% 以后我们会看到，这一概念在热力学和统计物理中是十分重要和有用的。有关开系的问题将在第三章以后讨论，目前暂不考虑。
	% 有的作者将孤立系定义为与其它物体完全没有相互作用的系统，有的作者将孤立系定义为与其它物体既没有物质、也没有能量交换的系统。我们采用后一定义。它既包括与其它物体完全没有相互作用，也包括处在恒定外场(重力场、静电场、静磁场等)的情形。在与其它物体没有物质和能量交换的情形下，撤去系统的某种内部约束，不破坏系统的孤立性。我们以后会遇到这种情形。
\end{remark}

% 宏观物质可以用很少的量表征。这种特性源于：
% \textit{宏观测量与原子时间尺度相比极其缓慢，与原子空间尺度相比十分粗糙。}
% 宏观体系忽略了系统内部每个粒子的具体运动，正如Anderson所说：\textbf{\textit{More is different.}}
% 而热力学便是\textit{唯象}地描述多粒子行为的宏观理论。

\section{热力学第零定律：平衡态}

\begin{definition}
	{平衡态}{equilibrium state}
	在给定的外界条件和内部约束下，系统的宏观性质保持稳定，各部分之间也没有净热流、净物质流等耗散通量，这种状态称为平衡态(equilibrium state)。孤立的宏观系统在允许的弛豫过程充分进行后可以达到平衡。
\end{definition}

% \begin{remark}
% 	系统由初始状态达到平衡态所需的特征时间称为弛豫时间(relaxation time)。弛豫时间可以很短也可以很长。以通常条件下的气体为例。通过分子的频繁碰撞，气体在10⁻¹⁰s左右就可以在小区域内建立局域平衡，而整个气体的平衡则要通过诸如扩散、热传导等过程才能实现。浓度的均匀化在气体中可能需要几分钟，在固体中则可能需要数小时、数星期甚至更长的时间。
% \end{remark}

\begin{remark}
	在平衡态之下，系统的宏观性质虽然不随时间改变，但组成系统的大量微观粒子仍处在不断的运动之中，只是这些微观粒子运动的统计平均效果不变而已。
	% 因此热力学的平衡态是一种动的平衡，常称为热动平衡。
	% 系统宏观物理量的数值仍会发生或大或小的涨落(fluctuation)，这种涨落在适当的条件下可以观察到。不过，对于宏观的物质系统，在一般情况下涨落是极其微小可以忽略的。在热力学中我们将不考虑涨落，而认为平衡态下系统的宏观物理量具有确定的数值。
\end{remark}

\begin{remark}
	% 前面给出了孤立系统平衡态的定义。
	平衡态的概念不限于孤立系统。
	与单一热库接触的系统可以和热库达到热平衡。若外界维持温差，系统即使宏观量不随时间变化，也可能持续输运热量，处于非平衡稳态。
	% 我们以后会看到，对于处在各种条件下的系统，热力学用相应的热力学函数作为判据判定系统是否处在平衡态，并导出存在相互作用的两个系统达到热动平衡的平衡条件。
\end{remark}

% 现在讨论如何描述一个热力学系统的平衡态。
% 前面已经说过，在平衡态之下，系统各种宏观物理量都具有确定值。
% 热力学系统所处的平衡态就是由其宏观物理量的数值确定的。由于宏观量之间的内在联系，表现为数学上存在一定的函数关系，这些宏观量不可能全部独立地改变。
% 我们可以根据问题的性质和考虑的方便选择其中几个宏观量作为自变量。
% 这些自变量本身可以独立地改变，我们所研究的系统的其它宏观量又都可以表达为它们的函数。这些自变量就足以确定系统的平衡态，我们称它们为状态参量；其它的宏观变量既然可以表达为状态参量的函数，便称为状态函数。我们通过几个具体例子加以说明。
% 假设所研究的系统是具有固定质量的化学纯的气体。气体装在一个封闭的容器里，具有确定的体积和压强。如果对气体加热，容易发现，气体的体积由于封闭在容器内未有显著的改变，但压强却增加了。因此要描述该气体的状态至少需要体积和压强两个参量，这两个参量是可以独立改变的。体积V描述气体的几何性质，叫做几何参量；压强p描述气体的力学性质，叫做力学参量。对于液体和各向同性的固体，也可以用体积V和压强p作为几何参量和力学参量来描述它们的平衡态。对于非各向同性的固体，几何参量和力学参量是应变张量和应力张量。在本课程中我们限于讨论各向同性的固体。

\begin{definition}
	{热平衡}{thermal equilibrium}
	两个内部已平衡的系统经导热壁接触，若没有净热量从一方流向另一方，它们就处于热平衡(thermal equilibrium)。
\end{definition}

\begin{theorem}{热力学第零定律：热平衡定律}{thermal equilibrium}
	系统A和B热平衡且A和C热平衡$\implies$ B和C热平衡。
\end{theorem}

\begin{corollary}
	互为热平衡的系统有一共同的物理性质，称为温度$T$，单位为 \unit{K}。
\end{corollary}

\begin{definition}
	{状态参量}{state parameter}
	为了描述一个热力学系统的平衡态，我们需要确定一组独立的宏观物理量，称为状态参量(state parameter)。
\end{definition}

% \begin{corollary}
% 	% 除补充说明外，
% 	系统的内能$U$、体积$V$和摩尔数$n$构成一组状态参量。
% \end{corollary}

\begin{remark}
	平衡态下很多物理量都具有确定值，但并不都是独立的，即存在一定的函数关系。
\end{remark}

\begin{definition}{状态方程}{state equation}
	状态方程(state equation)是温度$T$与其它状态参量间的关系。
\end{definition}

\begin{example}
	{理想气体状态方程}{ideal gas}
	历史上关于气体的实验\footnote{Boyle, Charles, Gay-Lussac, Clapeyron, Avogadro, Dalton, etc.}已经给出，在\textit{温度不太低、压强不太高}的条件下，气体的状态方程近似地有非常简单的规律。
	我们称满足这些规律的气体为理想气体(ideal gas)。

	在国际单位制下，
	理想气体的压强$p$、体积$V$、温度$T$和摩尔数$n$满足
	\footnote{摩尔数即物质的量(amount of substance)，与粒子数$N$成正比，单位为 \unit{mol}，即$n := N/\NA$，其中$\NA = \qty{6.02214076e23}{\per\mole}$是Avogadro常数。}
	\begin{equation}
		\label{eqn:pV=nRT}
		pV = nRT.
	\end{equation}
	其中$R = \qty{8.314}{\joule\per\kelvin\per\mole}$是理想气体常数。
\end{example}

\begin{example}
	{van der Waals气体状态方程}{van der Waals gas}
	van der Waals气体考虑了分子间的相互作用和分子体积：
	\begin{equation}
		\label{eqn:vdW}
		\biggp{p + a\frac{n^2}{V^2}}(V - nb) = nRT.
	\end{equation}
	% 其中$a$与分子间的相互作用有关，$b$与分子体积有关。
\end{example}

已经处于平衡态的热力学系统在外界影响下，会过渡到新的平衡态，在这个过程中的任一时刻，系统的状态都不是平衡态。
但是如果这个过程中的变化速度足够慢，每一瞬时都可以无限接近平衡态，就可以用平衡态处理整个过程。

\begin{definition}{准静态过程}{quasistatic process}
	准静态过程(quasi-static process)指每一瞬时，系统状态都无限接近平衡态的过程。
\end{definition}

\begin{corollary}
	在无摩擦等耗散机械作用的准静态过程中，外界对系统所做的功为
	\footnote{功的微元$\vd W$依赖于过程，一般不能写成一个功函数$W$的微分形式。}
	\begin{equation}
		\label{eqn:dW=Ydy}
		\vd W = \sum_iY_i\d y_i,
	\end{equation}
	其中$(Y_i,y_i)$分别是广义力和广义坐标，如压强--体积$(-p,V)$、表面张力--面积$(\sigma,A)$，以及外磁场--总磁矩$(B,m)$等。
	一般只考虑$\vd W = -p\d V$。
\end{corollary}

\section{热力学第一定律：能量守恒}

系统的能量包括内能$U$和整体运动能量。对于封闭系统，能量交换有功$W$和热量$Q$两种方式。

\begin{theorem}{热力学第一定律(能量守恒定律)}{first law of thermodynamics}
	对于宏观动能和外势能不变的封闭系统，内能增量$\d U$等于外界对它所做的功$\vd W$与外界向它传递的热量$\vd Q$的和：
	\begin{equation}
		\label{eqn:dU=dW+dQ}
		\d U = \vd W + \vd Q.
	\end{equation}
	% 由于功$\vd W$和热量$\vd Q$依赖于过程，因此不能写成微分形式。
\end{theorem}

\begin{remark}
	如果系统是\textbf{绝热}（$\vd Q \equiv 0$）的，我们便可以用机械功$\vd W$测量内能的变化$\Delta U$，通过指定基准态的内能$U_0$就可以得出任意状态的内能$U$。进而我们可以测量导热系统的传热$\vd Q$。
\end{remark}

% \begin{example}
% 	{理想气体的绝热过程}{}
% 	理想气体的内能
% 	\begin{equation}
% 		U=cnRT.
% 	\end{equation}
% 	其中$c$与理想气体粒子的自由度有关。由\eqref{eqn:pV=nRT}，$U=cpV$，继而
% 	\[
% 		\d U=c\d(pV)=c(p\d V+V\d p).
% 	\]
% 	再由式\eqref{eqn:dU=dW+dQ}，$\d U=-p\d V+0$，可得 
% 	\begin{equation}
% 		pV^\gamma=\const,\quad\gamma:=1+c^{-1}.
% 	\end{equation}
% 	绝热过程所做的功
% 	\[
% 		W=\frac{p_{\mathrm B}V_{\mathrm B}-p_{\mathrm A}V_{\mathrm A}}{1-\gamma}.
% 	\]
% \end{example}

\begin{definition}{热容}{heat capacity}
	定义热容(heat capacity)是物质在单位温度变化下所吸收或放出的热量：
	\begin{equation}
		\label{eqn:C general}
		C := \lim_{\Delta T\to0}\frac{\Delta Q}{\Delta T}.
	\end{equation}
	% 比热容(specific \textasciitilde)是单位质量的热容。
	% 摩尔热容$c:=C/n$是单位摩尔的热容。
	显然，热容与过程相关。
	以下取物质的量固定、仅有体积功的系统。等容过程中$\vd W = 0$，故等容热容可以写成
	\footnote{对于不同的过程，内能可能是不同状态参量的函数，比如$U = U(p,T)$和$U = U(V,T)$，尽管二者的偏导都写作$\p U/\p T$，但却是不同的函数形式。
	因此有必要将不变的状态参数写在偏导的脚标上。}
	\begin{equation}
		\label{eqn:C_V}
		C_V = \pvat{U}{T}{V}.
		\tag{\ref{eqn:C general}a}
	\end{equation}
	等压过程中$p = \const$，$\vd Q = \d U + p\d V$，故等压热容可以写成
	\begin{equation}
		\label{eqn:C_p}
		C_p = \pvat{U}{T}{p} + p\pvat{V}{T}{p}.
		\tag{\ref{eqn:C general}b}
	\end{equation}
\end{definition}

\begin{definition}
	{响应函数}{response function}
	响应函数描述一个物理量随另一个物理量的变化，通常可以直接测量。除了热容$C$，本篇采用如下热膨胀系数$\alpha$、压强系数$\beta$和等温压缩系数$\kappa_T$：
	\begin{subequations}
		\begin{align}
			\alpha &:= \frac1V\pvat{V}{T}{p},\\
			\beta &:= -\frac1p\pvat{p}{T}{V},\\
			\kappa_T &:= -\frac1V\pvat{V}{p}{T}.
		\end{align}
	\end{subequations}
\end{definition}

\begin{corollary}
	由$\d V/V = \alpha\d T - \kappa_T\d p$，在等容条件下有$(\p p/\p T)_V = \alpha/\kappa_T$，故
	\begin{equation}
		\alpha = -p\beta\kappa_T.
	\end{equation}
\end{corollary}

% \paragraph{内能标准全微分式}

% 将$U$全微分式中各变量微分前的系数用可测量表达出来。

\begin{theorem}
	{偏导关系式}{partial derivative relation}
	在推导不同的偏导数时，通常需要用到链式法则和隐函数定理：
	% 可以用 Jacobi 行列式证明。
	\begin{subequations}
		\begin{align}
			\label{eqn:XY->YX}
			\pvat{X}{Y}{Z} &= \pvat{Y}{X}{Z}^{-1};\\
			\label{eqn:XY*YZ*ZX=-1}
			\pvat{X}{Y}{Z} &= -\fracdisp{\pvat{Z}{Y}{X}}{\pvat{Z}{X}{Y}} \iff \pvat{X}{Y}{Z}\pvat{Y}{Z}{X}\pvat{Z}{X}{Y} = -1;\\
			\label{eqn:XY->XW/YW}
			\pvat{X}{Y}{Z} &= \fracdisp{\pvat{X}{W}{Z}}{\pvat{Y}{W}{Z}};\\
			\label{eqn:XY->XY+XW*WY}
			\pvat{X}{Y}{Z} &= \pvat{X}{Y}{W} + \pvat{X}{W}{Y}\pvat{W}{Y}{Z}.
		\end{align}
	\end{subequations}
\end{theorem}

\begin{example}
	{Joule自由膨胀实验}{Joule free expansion experiment}
	Joule在1845年用自由膨胀实验研究气体的内能。
	\begin{notecaptionblock}
		\tikzfigure{joule-expansion}
		\captionof{figure}{Joule的自由膨胀实验}
		\label{fig:Joule free expansion}
	\end{notecaptionblock}
	如\figref{fig:Joule free expansion} 所示，浸在水中的导热容器被活门隔开，一半填充压缩气体，另一半为真空。打开阀门让气体充满整个容器，实验在测量精度内未发现水温变化。

	向真空膨胀时$W = 0$。在理想气体极限，气体自由膨胀后和水重新平衡，二者的温度与初态相同，气体的净吸热为零。此时$(\p T/\p V)_U = 0$，利用$\d U = C_V\d T + (\p U/\p V)_T\d V$可得
	\[
		\pvat{U}{V}{T} = -\pvat{U}{T}{V}\pvat{T}{V}{U} = 0.
	\]
	即给定物质的量时理想气体的内能只取决于温度，称为Joule定律。
	\footnote{
		水和容器的总热容远大于气体的热容，实际气体膨胀引起的微小温度变化很难分辨。
		1852年，Joule和Thomson设计了节流膨胀实验来测量实际气体降压时的温度变化。
		尽管如此，Joule定律在低压极限下仍然适用。
		% 我们把严格遵循Boyle定律、Avogadro定律和Joule定律的气体称为理想气体。
	}
\end{example}

\begin{example}
	{理想气体的内能和比热}{ideal gas U and C}
	理想气体遵循Joule定律，$U = U(T)$，可以将等容热容写成导数：
	\begin{equation}
		\label{eqn:C_V->U}
		C_V = \dv UT \implies U = \int_{T_0}^TC_V\d T + U_0.
	\end{equation}
	在一定温度范围内，若理想气体的等容摩尔热容$C_V/n = cR$为常数，并取能量零点使$U_0 = cnRT_0$，则
	\begin{equation}
		U = cnRT.
	\end{equation}
	其中$c$与理想气体粒子的自由度有关。

	利用\eqref{eqn:XY->XY+XW*WY}，等容热容和等压热容的差
	\begin{align}
		\notag
		C_p - C_V &= \pvat{U}{T}{p} + p\pvat{V}{T}{p} - \pvat{U}{T}{V}\\
		\notag
		&= \cancel{\pvat{U}{T}{V}} + \pvat{U}{V}{T}\pvat{V}{T}{p} + p\pvat{V}{T}{p} - \cancel{\pvat{U}{T}{V}}\\
		&= \biggb{\pvat{U}{V}{T} + p}\pvat{V}{T}{p}\\
		&= (0 + p) \cdot \frac{nR}p = nR.
	\end{align}
	定义热容比
	\begin{equation}
		\gamma := \frac{C_p}{C_V} = \frac{c + 1}{c}.
	\end{equation}
	% 由此便可将$C_p,C_V$用$nR,\gamma$表示出来。
\end{example}

\begin{example}
	{理想气体的四种过程}{}
	理想气体有等容(isochoric, $\d V \equiv 0$)、等压(isobaric, $\d p \equiv 0$)、等温(isothermal, $\d T \equiv 0$)和绝热(adiabatic, $\vd Q \equiv 0$)等不同准静态过程。
	% 对于前三个过程，可简单利用\eqref{eqn:pV=nRT}得到关系。
	不同过程在$p\vs V$图中对应的路径如\figref{fig:ideal p-V} 所示，
	% 利用
	% \[
	% 	W=-\int_{V_\text A}^{V_\text B}p\d V,
	% \]
	% 计算功$W$。若方便，可用\eqref{eqn:C_V->U}计算内能变化$\Delta U$，再利用\eqref{eqn:dU=dW+dQ}计算热$Q=\Delta U-W$。
	\begin{notecaptionblock}
		\tikzfigure{ideal-gas-processes}
		\captionof{figure}{理想气体的$p\vs V$图}
		\label{fig:ideal p-V}
	\end{notecaptionblock}
	\tcblower
	下面研究理想气体的可逆绝热过程。
	由于$\vd Q = 0$，可得
	\[
		\d U = 0 + \vd W = -p\d V,
	\]
	另一方面
	\[
		\d U = C_V\d T
	\]
	此外对理想气体状态方程\eqref{eqn:pV=nRT}两边求全微分可得
	\[
		p\d V + V\d p = nR\d T = (C_p - C_V)\d T.
	\]
	结合以上三式，消去$\d T$。在$\gamma$可视为常数的温区内积分，得到
	\begin{equation}
		\label{eqn:adiabatic}
		\gamma p\d V + V\d p = 0 \implies pV^\gamma = \const.
	\end{equation}
	由于$\gamma > 1$，故$p\vs V$图中的绝热线会比等温线下降得更快，如\figref{fig:ideal p-V} 所示。
\end{example}

\section{热力学第二定律：熵}

热力学第一定律指出各种形式的能量在相互转化的过程中必须满足能量守恒定律，并没有限制过程进行的方向。
但在现实中，涉及热现象的过程都具有方向性。

给定一个初态未平衡的孤立系统，经过一段时间的演化会达到平衡态，这个过程未受环境影响，称为自发过程。
热量在有温差的物体间传递、气体自由膨胀等宏观弛豫过程具有方向性。达到平衡后，宏观系统在通常观测时间内自发回到原先非平衡态的概率极小。
比如热量可以自发地从高温物体传导到低温物体，但其逆过程却不可能自发发生：热量无法自发地从低温物体传到高温物体。

\begin{theorem}
	{热力学第二定律(Clausius表述)}{second law of thermodynamics (Clausius)}
	不可能把热量从低温物体传到高温物体，而不引起其它变化。
\end{theorem}

\begin{corollary}
	热量从低温物体传到高温物体，必须伴随其它变化，例如做功或热量的转移。
\end{corollary}

\begin{definition}
	{可逆过程}{reversible process}
	给定一个过程，若其逆过程可以使系统和外界都回到原来的状态而无其他变化，则称该过程为可逆过程(reversible process)。
\end{definition}

\begin{corollary}
	准静态过程若同时消除了摩擦、有限温差传热等熵产生机制，就可以可逆地进行。
	% 与热现象有关的实际过程都是不可逆过程。
\end{corollary}

\subsection{Carnot循环}

热机(heat engine)是实现热和功之间转换的一种系统。
Carnot在1824年研究热机时指出
对于在两个热源之间循环运行的热机，工作物质从高温热源$T_\hi$吸取热量，向低温热源$T_\lo$放出热量，两者的差额给出净输出功。

\begin{example}
	{理想气体的Carnot循环}{Carnot cycle of ideal gas}
	以理想气体为工作物质的Carnot热机，工作过程为四个准静态过程的循环：
	\footnote{约定记号：箭头上的量在过程中$\equiv 0$。
	% 可分别代表等容($\d V\equiv 0$)、等压($\d p\equiv 0$)、等温($\d T\equiv 0$)和绝热($\vd Q\equiv 0$)。
	}
	\[
		\text A(V_\text A,T_\hi)
		\overset{\d T} \longrightarrow
		\text B(V_\text B,T_\hi)
		\overset{\vd Q} \longrightarrow
		\text C(V_\text C,T_\lo)
		\overset{\d T} \longrightarrow
		\text D(V_\text D,T_\lo)
		\overset{\vd Q} \longrightarrow
		\text A(V_\text A,T_\hi).
	\]
	称为Carnot循环(carnot cycle)，如\figref{fig:Carnot cycle} 所示。\footnote{实线只是示意图，虚线才是正确的形状，单靠坐标伸缩是不可能变成实线那么好看的。}
	\begin{notecaptionblock}
		\tikzfigure{carnot-cycle}
		\captionof{figure}{Carnot循环}
		\label{fig:Carnot cycle}
	\end{notecaptionblock}
	取$Q_\hi,Q_\lo > 0$分别表示吸热量和放热量的大小。设两条绝热支所在温区内$C_V = cnR$为常数，则四步过程给出
	\begin{alignat*}{2}
		\text A \to \text B~&\text{等温膨胀(吸热)：} & Q_\hi &= nRT_\hi\ln\biggp{\frac{V_\text B}{V_\text A}};\\
		\text B \to \text C~&\text{绝热膨胀(对外做功)：} & \frac{V_\text C}{V_\text B} &= \biggp{\frac{T_\hi}{T_\lo}}^c;\\
		\text C \to \text D~&\text{等温压缩(放热)：} & Q_\lo &= nRT_\lo\ln\biggp{\frac{V_\text C}{V_\text D}};\\
		\text D \to \text A~&\text{绝热压缩(外界对气体做功)：} & \frac{V_\text D}{V_\text A} &= \biggp{\frac{T_\hi}{T_\lo}}^c.
	\end{alignat*}
	整个循环过程完成后，气体回到原来状态，对外所作的净功为
	\begin{equation}
		W_{\rm out} = Q_\hi - Q_\lo = nR(T_\hi - T_\lo)\ln\biggp{\frac{V_\text B}{V_\text A}}.
	\end{equation}
	热功转化的效率
	\begin{equation}
		\eta := \frac{W_{\rm out}}{Q_\hi} = 1 - \frac{T_\lo}{T_\hi} < 1
	\end{equation}
	在第一定律的符号约定下，一个循环中外界对工质做的净功为$W = -W_{\rm out}$。
\end{example}

\begin{corollary}
	Carnot循环是可逆过程。
\end{corollary}

\begin{theorem}{热力学第二定律(Kelvin表述)}{second law of thermodynamics (Kelvin)}
	不可能从单一热源吸热，使之完全变成有用功，而不引起其它变化。
\end{theorem}

\begin{corollary}
	两种表述是等价的。
\end{corollary}

\begin{proof}
	考虑Carnot循环，热机从高温热源$T_\hi$吸取热量$Q_\hi$、向低温热源$T_\lo$放出热量$Q_\lo$，对外做功$W_{\rm out} = Q_\hi - Q_\lo$。
	如果Clausius表述不成立，则可以将热量$Q_\lo$从$T_\lo$传到$T_\hi$而不引起其它变化。这相当于热机从单一热源$T_\hi$吸取热量$Q_\hi - Q_\lo$，并将其完全转化为有用功，与Kelvin表述矛盾。

	反之，如果Kelvin表述不成立，则可以从单一热源$T_\lo$吸取热量$Q_\lo$，并将其完全转化为有用功$W_{\rm out} = Q_\lo$，而不引起其它变化。将此功完全转化为$T_\hi$的热量，就相当于热量从$T_\lo$传到$T_\hi$而没有引起其他变化，这与Clausius表述矛盾。
\end{proof}

\begin{theorem}{Carnot定理}{Carnot's Theorem}
	在相同高、低温热源之间工作的热机的效率
	\begin{equation}
		\eta = 1 - \frac{Q_\lo}{Q_\hi} \leq 1 - \frac{T_\lo}{T_\hi}.
	\end{equation}
	取等号当且仅当热机是可逆的。
\end{theorem}

\begin{proof}
	令可逆热机A和待比较的热机B从高温热源各吸收$Q_\hi$，对外输出功分别为$W_{\rm A},W_{\rm B}$。若$W_{\rm B} > W_{\rm A}$，令A反向运行，两机合起来只从低温热源取热$W_{\rm B} - W_{\rm A}$，并将其全部转化为功，与Kelvin表述矛盾。因此$\eta_{\rm B} \leq \eta_{\rm A}$。
	若B也可逆，交换A、B即得反向不等式，两者效率相等。反之，效率相等时，反向运行A恰好复原B对两个热源及功源的全部影响，符合可逆过程的定义。
\end{proof}

\begin{corollary}
	可逆热机效率只与热源温度有关，与工作物质无关。
\end{corollary}

\begin{definition}
	{热力学温标}{thermodynamic temperature scale}
	借助Carnot热机，可通过测量热源吸收/放出的热量的比值来确定温度
	\begin{equation}
		\label{eq:sm-absolute-temperature-ratio}
		\frac{Q_\lo}{Q_\hi} = \frac{T_\lo}{T_\hi}.
	\end{equation}
	这一比值与工作物质无关，由此建立的绝对温标称为热力学温标(thermodynamic temperature scale)。历史上以指定同位素组成的水的三相点为参考，赋予其温度$\qty{273.16}{\kelvin}$。现行国际单位制通过固定Boltzmann常数$\kB$的数值定义开尔文。
\end{definition}

% \begin{remark}
% 	为了准确测定温度，我们不止需要指定温度的特殊点，还需要指定测温物质和测温属性。例如，单纯利用物体的热膨胀测温，那么水银温度计和酒精温度计将会给出不同的温度结果。而热力学温标与测温物质无关的特性，使得它在物理中被广泛运用。
% \end{remark}

\subsection{熵}

\begin{theorem}{Clausius不等式}{Clausius inequality}
	封闭系统在一个循环中从温度为$T_\reservoir$的热源吸收热量$\vd Q$。各次换热可对应不同的热源温度，且满足
	\begin{equation}
		\oint\frac{\vd Q}{T_\reservoir} \leqslant 0.
	\end{equation}
	当且仅当为可逆循环时取等号。
\end{theorem}

\begin{proof}
	设系统完成一个循环，从热源$T_i$吸收的净热量为$Q_i$，$Q_i$可以为负。为每个热源配一台可逆辅助热机，使其向该热源返还$Q_i$。由可逆热机的热量比，它同时从一个公共参考热源$T_0$吸收$T_0 Q_i/T_i$。各台热机可根据$Q_i$的符号正向或反向运行。
	这样，系统与辅助热机完成循环后，只同参考热源有净换热。若从参考热源吸收的净热量为正，该热量就全部变成了功，与Kelvin表述矛盾。因此
	\[
		T_0\sum_i\frac{Q_i}{T_i} \leq 0 .
	\]
	若原循环可逆，反向运行使各$Q_i$变号，因此必须取等号。反之，取等号时，辅助热机恰好复原了原循环给各热源和功源带来的变化，原循环可逆。让热源温度连续变化，求和便成为积分。
\end{proof}

\begin{corollary}
	考虑任意经过$\text A \to \text B$的两个可逆过程$\text R,\text R'$，并记$\text R'$的逆过程为$-\text R'$，则$\text R + (-\text R') \equiv \text R - \text R'$构成一个可逆循环，由Clausius不等式：
	\[
		\oint_{\text R-\text R'}\frac{\vd Q}T = \int_\text R\frac{\vd Q}T + \int_{-\text R'}\frac{\vd Q}T = 0 \implies
		\int_\text R\frac{\vd Q}T = \int_{\text R'}\frac{\vd Q}T,
	\]
	因此对于经过$\text A \to \text B$的任意可逆过程，$\vd Q/T$的积分都是相同的，与过程的具体路径无关，由此可写成微分的形式。
\end{corollary}

\begin{definition}
	{熵}{entropy}
	对于可逆过程，可定义微分
	\begin{equation}
		\label{eqn:dS=dQ/T}
		\d S := \frac{\vd Q}T.
	\end{equation}
	称$S$为熵(entropy)。
\end{definition}

\begin{corollary}
	设封闭系统从平衡态A经过任意过程到达平衡态B。用可逆过程从B返回A，应用Clausius不等式得到
	\begin{equation}
		\label{eq:sm-entropy-clausius-inequality}
		S_\mathrm{B} - S_\mathrm{A} \geq \int_\mathrm{A}^\mathrm{B}\frac{\vd Q}{T_\reservoir} .
	\end{equation}
	有熵产生时取严格不等号。右边用每次换热时的热源温度，因而也适用于中途系统温度不均匀的过程。
\end{corollary}

\begin{theorem}
	{熵增加原理}{principle of entropy increase}
	在绝热条件下，封闭系统在可逆过程中$\Delta S = 0$，有熵产生时$\Delta S > 0$。
\end{theorem}

\begin{corollary}
	孤立系统在给定总内能、体积和物质的量时，平衡态的熵最大。对可加系统，在固定物质的量下，平衡熵$S(U,V)$是凹函数。在两相共存区，改变两相比例可使熵沿某些方向线性变化。
	% \[
	% 	\pv[2]SU\pv[2]SV>\biggp{\pw SUV}^2>0.
	% \]
\end{corollary}

% 熵是热运动混乱程度的量度。

% \begin{example}{熵的计算}{Calculating Entropy}
% 	将质量相同而温度分别为$T_1$和$T_2$的两杯水在等压下绝热的混合，求熵变。

% 	终态温度$T=(T_1+T_2)/2$，第一杯水的熵变为
% 	\[
% 		\Delta S_1=\int_{T_1}^T\frac{C_p\d T}T=C_p\ln\frac{T_1+T_2}{2T_1},
% 	\]
% 	第二杯水的熵变$\Delta S_2$同理可求，
% 	总熵增
% 	\[
% 		\Delta S_1+\Delta S_2=C_p\ln\frac{(T_1+T_2)^2}{4T_1T_2}\geqslant 0.
% 	\]
% 	取等号当且仅当$T_1=T_2$。
% \end{example}

\begin{theorem}
	{热力学基本微分方程}{}
	对于物质的量固定、仅作体积功的简单系统，沿可逆路径有$\vd Q = T\d S$和$\vd W = -p\d V$。代入第一定律，得到相邻平衡态之间的状态函数微分：
	\begin{equation}
		\label{eqn:dU=TdS-pdV}
		\d U = T\d S - p\d V.
	\end{equation}
\end{theorem}

\begin{corollary}
	根据全微分与偏导数的关系，温度和压强可以写成偏导形式：
	\begin{subequations}
		\begin{align}
			T &= \pvat{U}{S}{V},\\
			p &= -\pvat{U}{V}{S}.
		\end{align}
	\end{subequations}
	在熵表象下，考虑$S$的全微分，可由\eqref{eqn:dU=TdS-pdV}直接得出
	\begin{equation}
		\d S = \frac1T\d U + \frac pT\d V.
	\end{equation}
	以及$S$的偏导数
	\begin{subequations}
		\begin{align}
			\pvat{S}{U}{V} &= \frac1T,\\
			\pvat{S}{V}{U} &= \frac pT.
		\end{align}
	\end{subequations}
\end{corollary}

% \begin{theorem}
% 	{熵函数的凹凸性}{concavity of entropy function}
% 	$S=S(U,V)$是严格上凸的(strictly concave)。
% \end{theorem}

% \begin{proof}
% 	固定$V$，对于内能为$2U$的系统，其两个子系统内能$U_\pm=U\pm\Delta U$。不均匀的子系统熵更低，即
% 	\[
% 		S(U+\Delta U)+S(U-\Delta U)<S(2U)=2S(U)
% 		\implies\pv[2]SU<0,
% 	\]
% \end{proof}

\begin{corollary}
	在固定体积和物质的量的严格稳定单相区，熵对内能的二阶导数为负。因此$1/T = (\p S/\p U)_V$随$U$增加而减小。在正温区，增加内能会使温度升高。
\end{corollary}

\begin{example}
	{温度$T$的热力学意义}{}
	将孤立系统用固定的导热壁隔开，两部分只能相互交换热量，子系统的内能满足
	\[
		U_1 + U_2 = \const \implies \d U_1 = -\d U_2.
	\]
	平衡时熵达到极值，有
	\[
		\d S = \d S_1 + \d S_2 = \pv{S_1}{U_1}\d U_1 + \pv{S_2}{U_2}\d U_2 = \biggp{\frac1{T_1} - \frac1{T_2}}\d U_1 = 0 \implies T_1 = T_2.
	\]
	故平衡时子系统温度相等；
	若两部分均为热容为正的单相系统，初始时$0 < T_1^0 < T_2^0$，则由能量守恒和内能对温度的单调性，平衡温度$T$满足
	\[
		T_1^0 < T < T_2^0,
	\]
	这说明\textit{热量总是从高温流向低温}，这符合现实中温度的特性。
\end{example}

\begin{example}
	{压强$p$的热力学意义}{}
	系统由导热
	% \footnote{熵表象下的绝热活塞讨论比较微妙。}
	活塞（$U,V$可变）隔开，平衡时
	\[
		\d S = \d S_1 + \d S_2 = \biggp{\frac1{T_1} - \frac1{T_2}}\d U_1 + \biggp{\frac{p_1}{T_1} - \frac{p_2}{T_2}}\d V_1 = 0.
	\]
	能量和体积可独立变化，故$T_1 = T_2$，进而有$p_1 = p_2$。同理，可以进而推出\textit{压强大的一侧会推动活塞}。
	% 换成导热半透膜（$V$恒定，$U,N$可变）隔开，可以得到$T_1=T_2,N_1=N_2$的结论，说明\textit{物质总是由化学势高流向化学势低的一边}。
\end{example}

\begin{example}
	{理想气体的熵}{entropy of ideal gas}
	考虑理想气体的微分
	\[
		\d S = \frac1T\d U + \frac pT\d V = \frac{C_V}T\d T + \frac{nR}V\d V.
	\]
	采用$C_V = cnR$，可得
	\begin{equation}
		\label{eq:sm-ideal-gas-entropy}
		\d S = nR\biggp{c\frac{\d T}T + \frac{\d V}V} \implies S = nR\ln\biggb{\biggp{\frac{T}{T_0}}^c\biggp{\frac{V}{V_0}}} + S_0.
	\end{equation}
	因此，对于可逆绝热过程，$\d S = 0$，从而$T^cV = \const$。
\end{example}

\paragraph{用过程求熵}

因为熵的变化只与始末态有关，可以选取连接这两个平衡态的可逆路径，%我们可以让这个过程有比较良好的性质：各子系统无限接近平衡。由于
%$$\d U=\vd Q+\vd W+\vd W_{\mathrm r}.$$
%其中$\vd Q$为传热，$\vd W=-P\d V$为功，$\vd W_{\mathrm r}$为耗散，则熵的变化
%$$\d S=\frac1T\d U+\frac PT\d V=\frac{\vd Q+\vd W_{\mathrm r}}T.$$
%因为在准静态过程中，\textit{只有热和耗散导致了熵的增加}。我们可以构造特定的过程，
用沿途的可逆吸热计算熵变
\begin{equation}
	\label{eq:sm-reversible-entropy-change}
	\Delta S = \int\frac{\vd Q_\rev}{T}.
\end{equation}

\begin{example}
	{理想气体自由膨胀}{}
	理想气体从$V$自由(绝热)膨胀到$\lambda V$，其熵增为
	\begin{equation}
		\label{eq:sm-free-expansion-entropy}
		\Delta S = nR\ln\lambda.
	\end{equation}
	显然$\Delta S > 0$，然而实际上$Q \equiv 0$。因此自由膨胀是非准静态的，更不可逆。
	\tcblower
	自由膨胀的初末态A、C具有相同温度。膨胀中途气体不均匀，无法在平衡态的$p\vs V$图上描出实际路径。
	为了求其熵变，可以构造先$\text A \to \text B$可逆绝热膨胀，后$\text B \to \text C$可逆等容加热的路径
	\[
		\text A(V,p_\text A)\overset{\vd Q} \longrightarrow \text B(\lambda V,p_\text B)\overset{\d V} \longrightarrow \text C(\lambda V,p_\text C).
	\]
	\begin{notecaptionblock}
		\tikzfigure{expansion-paths}
		\captionof{figure}{连接A、C的可逆等温路径和经过B的另一条可逆路径}
	\end{notecaptionblock}
	$\text A \to \text B$绝热，$T^cV = \const.$因此$T_\text A^c = \lambda T_\text B^c.$\\[1ex]
	$\text B \to \text C$等容，$\vd Q = \d U = cnR\d T$，故
	\[
		\Delta S = \int_\text B^\text C\frac{\vd Q_\rev}T = cnR\int_{T_\mathrm{B}}^{T_\mathrm{C}}\frac{\d T}T = nR\ln\lambda.
	\]
	与$\text A \to \text C$等温过程的熵变相同。
\end{example}


\section{热力学第三定律：绝对零度}

相较于前几个基本定律，
热力学第三定律的出现得非常晚。
本节适合在学习完\chapref{chap:simple system}后阅读。

\subsectionstar{低温的获得}

% 物理中有各种能量尺度。 在一定的温度下，相应能量尺度（∼ kBT）的现象会显现出来。
为了探索更低能量尺度上的物理，人们需要越来越低的温度。早期获得低温的一条重要途径是气体液化。1877年观察到氧液化，1883年制得可保持的液氧和液氮，1898年和1908年分别制得液氢和液氦。液化气体在给定压强下沸腾时可作为低温热浴，用来冷却与它热接触的样品。
% 液化气体的关键在于降低气体分子的动能，这样由于分子间的吸引力，可将气体液化。降低动能有各种不同的方法，但基本上都是使气体绝热膨胀。实现绝热膨胀的方法不同，对温度的改变也不同。气体液化后再降温会比较困难，但液化后的气体有很大的潜热，可用于保持恒温。

\begin{example}
	{Joule--Thomson节流膨胀实验}{}
	% 鉴于1843年，Joule的自由膨胀实验不够精确，1852年Joule和Thomson设计了一个节流膨胀实验来观察实际气体在膨胀时所发生的温度变化。
	如\figref{fig:Joule throttle expansion} 所示，在一个圆形绝热筒的中部，置有一个刚性的多孔塞，
	实验时，将压力和温度恒定为$p_1$和$T_1$的某种气体，连续地压过多孔塞，使气体在多孔塞右边的压力恒定为$p_2$，且$p_1 > p_2$。
	\footnote{由于多孔塞的孔很小，气体只能缓慢地从左侧进入右侧，并且可保持左室$p_1$部分和右室低压$p_2$的部分压力恒定不变。这种维持一定压力差的绝热膨胀过程叫做节流膨胀。}
	\begin{notecaptionblock}
		\tikzfigure{joule-thomson-expansion}
		\captionof{figure}{Joule--Thomson节流膨胀实验}
		\label{fig:Joule throttle expansion}
	\end{notecaptionblock}
	将1中所有的气体推到2后，
	\[
		\Delta U = U_2 - U_1 = -p_2V_2 + p_1V_1 + 0,
	\]
	即
	\begin{equation}
		\label{eq:sm-throttling-enthalpy}
		U_1 + p_1V_1 = U_2 + p_2V_2 \implies H_1 = H_2.
	\end{equation}
	因此节流膨胀过程焓变$\Delta H = 0$。
\end{example}

\begin{remark}
	节流时气体在多孔塞内存在压强梯度。上面的等焓关系联系入口和出口的平衡状态，条件是稳态绝热、无轴功，且两端的宏观动能和重力势能之差可以忽略。
\end{remark}

\begin{theorem}
	{Joule--Thomson效应}{}
	由\eqref{eqn:XY->XW/YW}和Maxwell关系\eqref{eqn:Maxwell relation G}：
	\begin{equation}
		\pvat{T}{p}{H} = -\Biggdivision{\pvat{H}{p}{T}}{\pvat{H}{T}{p}} = -\frac1{C_p}\biggb{T\pvat{S}{p}{T} + V} = \frac{V(T\alpha - 1)}{C_p}.
	\end{equation}
	当$T\alpha > 1$时，高压向低压的节流过程可以降低温度。
\end{theorem}
\begin{remark}
	理想气体$T\alpha = 1$，不能通过节流过程改变温度。
	但实际气体在一定条件下会有相当大的Joule--Thomson效应，Kamerlingh Onnes就是利用这种效应首次将He液化。
\end{remark}

% 若要获得尽可能低的温度，一种方法是让气体绝热膨胀；另一种是自由膨胀（所幸现实中的气体并不理想）。
\begin{example}
	{通过van der Waals气体实现降温}{}
	考虑van der Waals气体，其热膨胀系数
	\[
		\alpha = \frac1v\pvat{v}{T}{p} = \Bigp{\frac{Tv}{v - b} - \frac{2a(v - b)}{Rv^2}}^{-1}
	\]
	在力学稳定的单相分支上，$\alpha$的分母为正。代入节流制冷条件$T\alpha > 1$，得到
	\begin{equation}
		\label{eq:sm-van-der-waals-inversion-temperature}
		T < T_{\mathrm{inv}}(v),\qquad
		T_{\mathrm{inv}}(v) = \frac{2a}{bR}\biggp{1 - \frac bv}^2.
	\end{equation}
	等号对应反转点。沿反转曲线将$b/v = 1 - \sqrt{RbT/(2a)}$代回状态方程，可得
	\begin{equation}
		\label{eq:sm-van-der-waals-inversion-curve}
		p_{\mathrm{inv}}(T) = \frac a{b^2}\biggp{1 - \sqrt{\frac{RbT}{2a}}}\biggp{3\sqrt{\frac{RbT}{2a}} - 1}.
	\end{equation}
	低密度气体的反转温度趋于$2a/(Rb)$。在这一极限下，低于反转温度时降压制冷，高于反转温度时降压升温。
	\begin{notecaptionblock}
		\tikzfigure{van-der-waals-cooling}
		\captionof{figure}{van der Waals气体在正压下的节流反转曲线。右端点对应低密度极限。}
	\end{notecaptionblock}
\end{example}

\subsection{Nernst定理}

1906年，Nernst在研究各种化学反应在低温下的性质时，由于
\[
	\frac{\Delta G - \Delta H}T = -\Delta S
	% \implies\lim_{T\to 0}\frac{\Delta G}T-\lim_{T\to 0}\frac{\Delta H}T=-\lim_{T\to 0}(\Delta S)_T.
\]
仅有$\Delta H - \Delta G \to 0$还不能确定两条曲线在零温附近的斜率。Nernst热定理对其差值随温度趋零的速率给出了更强的限制。

\begin{theorem}
	{Nernst定理}{Nernst theorem}
	对处于平衡的纯物质，若在同一相内改变压强等外参量，避开改变基态简并结构的奇点，则Nernst热定理给出有限等温变化的低温极限
	\begin{equation}
		\lim_{T\to 0}(\Delta S)_T = 0.
	\end{equation}
\end{theorem}

\begin{corollary}
	热定理要求零温极限的熵与上述外参量无关。对于基态唯一的纯净完美晶体，Planck表述给出
	\begin{equation}
		\lim_{T\to 0}S(T,p) = 0.
	\end{equation}
	若基态有$g_0$重简并，在等权平衡态中有$S(0) = \kB\ln g_0$。若简并数随粒子数指数增长，单位物质的量也会有有限的剩余熵。
\end{corollary}

\begin{theorem}
	{热力学第三定律}{third law of thermodynamics}
	不可能通过有限的步骤使一个物体冷却到热力学温度的零度。
\end{theorem}

